\documentclass[../../main.tex]{subfiles}% 注意这里的文件路径不能用 ./main.tex ,否则用latexmk编译子文件会报错
\graphicspath{{\subfix{./image/}}{\subfix{../../image/}}} % 指定图片引用路径,后续可以直接使用图片文件名
% 如果图片存放在上级目录,则使用 ../../image/ 作为备用路径

% 例如:
% \begin{figure}[H]
% \centering
% \includegraphics[width=0.3\textwidth]{图.png}
% \caption{}
% \label{figure:图}
% \end{figure}
% 注意:上述\label{}一定要放在\caption{}之后,否则引用图片序号会只会显示??.

\begin{document}

\section{存在性与正则性结果}\label{section:ch06-存在正则}

\subsection{存在性}\label{subsection:ch06-存在性}

第一个存在性结果的证明是构造性的。事实上，我们通过研究逼近问题序列的解 $\{u_n\}$ 来展示 (P) 的弱解：
\begin{align*}
\begin{cases}
\partial_t u_n=\Delta\phi_n(u_n) & \text{在 }Q,\\
u_n=0 & \text{在 }\Gamma,\\
u_n(\cdot,0)=u_{0n} & \text{在 }\Omega,
\end{cases}
\qquad(Pn)
\end{align*}
（见\cref{section:ch06-逼近问题}）。此处 $\phi_n(z)\triangleq \phi(z)+\frac{z}{n}$（$z\in\mathbb{R}$），$u_{0n}$ 为初值测度 $u_0$ 的合适正则化（见\eqref{align:ch6-6.48}--\eqref{align:ch6-6.49a}）。

\begin{theorem}\label{theorem:ch06-6.4.1}
设 (A0)--(A2) 成立。则存在问题 (P) 的弱解，它作为逼近问题 (Pn) 解序列在 $L^\infty_{w^*}(0,T;\mathscr{R}_f(\Omega))$ 的弱$^*$拓扑下的极限点而得到。进而：
\begin{enumerate}[(i)]
\item 对几乎所有 $0<t_1<t_2<T$，
\begin{align}
u^\pm_s(\cdot,t_2)\leqslant u^\pm_s(\cdot,t_1)\leqslant u^\pm_{0s} \quad\text{在 }\mathscr{R}_f(\Omega). \label{align:ch6-6.33}
\end{align}
\item 若 $u_0\in L^1(\Omega)$，则 $u\in L^\infty(0,T;L^1(\Omega))$。
\end{enumerate}
\end{theorem}

由\cref{theorem:ch06-6.3.1}，为得到 (P) 的适定性，证明满足相容性条件的弱解存在性是重要的。这由如下结果保证。

\begin{theorem}\label{theorem:ch06-6.4.2}\label{proposition:ch06-6.4.2}
设 (A0)--(A2) 成立。则\cref{theorem:ch06-6.4.1}所给出的问题 (P) 的弱解满足相容性条件。
\end{theorem}

由\cref{theorem:ch06-6.3.1}、\cref{theorem:ch06-6.4.1}与\cref{theorem:ch06-6.4.2}，有如下结果。

\begin{theorem}\label{theorem:ch06-6.4.3}
设 (A0)--(A2) 成立。则对每个 $u_0\in\mathscr{R}_f(\Omega)$，存在问题 (P) 满足相容性条件的唯一弱解 $u$。进而\cref{theorem:ch06-6.4.1}的 (i)--(ii) 款成立。
\end{theorem}

\cref{theorem:ch06-6.4.1}证明中构造的 (P) 弱解有若干重要性质（特别地，由\cref{theorem:ch06-6.4.2}它们满足相容性条件）。这引出如下定义。

\begin{definition}\label{definition:ch06-构造解}
\cref{theorem:ch06-6.4.1}所给出的问题 (P) 弱解称为构造解。
\end{definition}

构造解的另一个有趣性质如下。

\begin{proposition}\label{proposition:ch06-6.4.4}
设 (A0)--(A2) 成立。则问题 (P) 的每个构造解都是熵解。
\end{proposition}

\subsection{正则性}\label{subsection:ch06-正则性}

熵解一个有趣的特征是其奇异部分既不能自发出现，也不能随时间增长：

\begin{theorem}\label{theorem:ch06-6.4.5}
设 (A0)--(A1) 成立。假设
\begin{align}
\text{存在 }L_0\geqslant0\text{ 使当 }|s|>L_0\text{ 时 }\phi'(s)>0. \label{align:ch6-6.34}
\end{align}
设 $u$ 为问题 (P) 的熵解。则对几乎所有 $t_1,t_2\in(0,T)$，$t_1<t_2$，不等式\eqref{align:ch6-6.33}成立。
\end{theorem}

\begin{proof}
设 $u$ 为问题 (P) 的熵解。先证对所有 $\rho\in C^1_c(\Omega)$，$\rho\geqslant0$，及 $h\in C^1([0,T])$，$h(T)=0$，有
\begin{align}
\int_0^T\langle u^\pm_s(\cdot,t),\rho\rangle h'(t)\,\mathrm{d}t\geqslant-h(0)\langle u^\pm_{0s},\rho\rangle. \label{align:ch6-6.35}
\end{align}
我们仅证 $+$ 号情形。因 $\phi$ 在 $[L_0,+\infty)$ 上严格增（见\eqref{align:ch6-6.34}），由标准逼近结果可在\eqref{align:ch6-6.11}中取 $F_\phi(y)=[y-m]_+$，$m>L_0$（这对应于在\eqref{align:ch6-6.9}中取 $f=\chi_{[\phi(m),+\infty)}$）。此时 $M^+_{F_\phi}=1$，$M^-_{F_\phi}=0$，从而得
\begin{align}
\begin{aligned}
&\iint_Q\{[u_r-m]_+\zeta_t-\chi_{[m,+\infty)}(u_r)\nabla\phi(u_r)\cdot\nabla\zeta\}\,\mathrm{d}x\,\mathrm{d}t\\
&\quad+\int_0^T\langle u^+_s(\cdot,t),\zeta_t(\cdot,t)\rangle\,\mathrm{d}t
\geqslant-\int_\Omega[u_{0r}-m]_+\zeta(x,0)\,\mathrm{d}x-\langle u^+_{0s},\zeta(\cdot,0)\rangle
\end{aligned} \label{align:ch6-6.36}
\end{align}
对每个 $\zeta\in C^1([0,T];C^1_c(\Omega))$，$\zeta\geqslant0$，$\zeta(\cdot,T)=0$ 在 $\Omega$ 上成立。在此不等式中取 $\zeta(x,t)=\rho(x)h(t)$，$\rho,h$ 如上，并令 $m\to\infty$，即得\eqref{align:ch6-6.35}。

现在固定 $t_1\in(0,T)$。对充分大的 $p\in\mathbb{N}$，置
\begin{align*}
h_p(t)\triangleq \chi_{[0,t_1]}(t)+p\left(t_1+\frac{1}{p}-t\right)\chi_{(t_1,t_1+\frac{1}{p}]}(t) \quad(t\in[0,T]).
\end{align*}
在\eqref{align:ch6-6.35}中取 $h=h_p$，对所有非负 $\rho\in C^1_c(\Omega)$，得
\begin{align}
p\int_{t_1}^{t_1+\frac{1}{p}}\langle u^\pm_s(\cdot,t),\rho\rangle\,\mathrm{d}t\leqslant\langle u^\pm_{0s},\rho\rangle. \label{align:ch6-6.37}
\end{align}
因由\cref{definition:ch06-熵解}，映射 $t\mapsto\langle u^\pm_s(\cdot,t),\rho\rangle\in L^\infty(0,T)$，由\cref{corollary:ch02-平均收敛}(i) 有
\begin{align*}
\lim_{p\to\infty}p\int_{t_1}^{t_1+\frac{1}{p}}\langle u^\pm_s(\cdot,t),\rho\rangle\,\mathrm{d}t
=\langle u^\pm_s(\cdot,t_1),\rho\rangle.
\end{align*}
在\eqref{align:ch6-6.37}中令 $p\to\infty$，对几乎所有 $t_1\in(0,T)$，得 $\langle u^\pm_s(\cdot,t_1),\rho\rangle\leqslant\langle u^\pm_{0s},\rho\rangle$，对所有 $\rho\in C^1_c(\Omega)$，$\rho\geqslant0$ 成立，故\eqref{align:ch6-6.33}中第二个不等式直接成立。第一个不等式类似地在\eqref{align:ch6-6.35}中取如下 $h=h_p$ 得证：
\begin{align*}
h_p(t)\triangleq p(t-t_1)\chi_{[t_1,t_1+\frac{1}{p}]}(t)
+\chi_{(t_1+\frac{1}{p},t_2)}(t)
+p\left(t_2+\frac{1}{p}-t\right)\chi_{[t_2,t_2+\frac{1}{p}]}(t)
\quad(t\in[0,T]).
\end{align*}
结论随之成立。
\end{proof}

\cref{theorem:ch06-6.4.5}可视为熵解的一个正则性结果。特别地，由\cref{proposition:ch06-6.4.4}，它适用于构造解。值得注意的是，(P) 的构造解还表现出进一步的正则化效应。为讨论这一点，陈述如下定义。

\begin{definition}\label{definition:ch06-RfL1正则}
问题 (P) 的弱解 $u$ 经历 $\mathscr{R}_f$--$L^1$ 正则化效应，若存在等待时间 $\tau\in[0,T)$，使对所有 $t\in(\tau,T)$ 有 $u(\cdot,t)\in L^1(\Omega)$。若 $\tau=0$，则发生瞬时 $\mathscr{R}_f$--$L^1$ 正则化。
\end{definition}

$\mathscr{R}_f$--$L^1$ 正则化效应发生的充分条件由如下命题给出，它依赖于 $\phi$ 趋于饱和值 $\phi_{\pm\infty}$ 的合适收敛速率（证明见文献~\cite{ref87} 命题 3.8）。

\begin{proposition}\label{proposition:ch6-瞬时正则}
设 (A0)--(A2) 成立。设存在 $s_\pm\in(0,1)$ 与 $M_\pm\in(0,\phi_{\pm\infty}]$，使
\begin{align}
\pm\phi_{\pm\infty}\mp\phi(y)\geqslant\frac{M_\pm}{(1+|y|)^{s_\pm}}
\quad\text{当 }\pm y\geqslant0. \label{align:ch6-6.38}
\end{align}
进而设 $u^\pm_0\in\mathscr{R}^+_{d,p}(\Omega)$，$p\in(1,\frac{2}{1+s_\pm})$，$u$ 为 (P) 满足相容性条件的唯一弱解。则 $u^\pm\in L^\infty(0,T;L^1(\Omega))$。
\end{proposition}

相反地，如下两个命题指出瞬时 $\mathscr{R}_f$--$L^1$ 正则化不发生的情形：或对合适的 $u_0\in\mathscr{R}_{d,2}(\Omega)$ 且 $\phi$ 严格增，或在 $\phi$ 在 $\pm\infty$ 附近为常数的最奇异情形下对任意构造解。证明见文献~\cite{ref87} 命题 3.9、3.10。

\begin{proposition}\label{proposition:ch06-6.4.7}
设 (A0)--(A2) 成立且 $\phi$ 严格增。则存在 $u_0\in\mathscr{R}^+_{d,2}(\Omega)$，$u_{0s}\neq0$，使对 (P) 满足相容性条件的唯一弱解，有
\begin{align}
\operatorname*{ess\,lim}_{t\to0^+}\|u_s(\cdot,t)-u_{0s}\|_{\mathscr{R}_f(\Omega)}=0. \label{align:ch6-6.39}
\end{align}
\end{proposition}

\begin{proposition}\label{proposition:ch06-6.4.8}
设 (A0)--(A2) 成立，$u$ 为 (P) 满足相容性条件的唯一弱解。设存在 $\pm c_\pm\geqslant0$ 与 $\pm d_\pm\geqslant0$，使 $\phi(y)=c_\pm$ 当 $\pm y\geqslant\pm d_\pm$。则对每个开子集 $\widetilde\Omega\Subset\Omega$，有
\begin{align}
\operatorname*{ess\,lim}_{t\to0^+}\|u^\pm_s(\cdot,t)-u^\pm_{0s}\|_{\mathscr{R}_f(\widetilde\Omega)}=0. \label{align:ch6-6.40}
\end{align}
\end{proposition}




\end{document}